\section*{Capítulo \ref{ch:g8:ice-on-earth} --- El hielo en la Tierra}

\begin{solution}{exo:g8:ice-on-earth:1}
(a) La zona de acumulación. (b) La zona de ablación. (c) El frente. (d) Una
\omterm{def:g8:ice-on-earth:glacial-landforms}{morrena} (frontal).
\end{solution}

\begin{solution}{exo:g8:ice-on-earth:2}
(a) Falso: un \omterm{def:g8:ice-on-earth:glacier}{glaciar} está hecho de \omterm{def:g4:clouds-rain-snow:precipitation}{nieve} caída sobre tierra firme y
transformada en hielo; el agua de \omterm{def:g2:water-on-earth:ocean}{mar} congelada es el \omterm{def:g8:ice-on-earth:sea-ice}{hielo marino}.
(b) Verdadero. (c) Falso: su hielo sigue fluyendo cuesta abajo; el frente
retrocede porque el hielo se funde más deprisa de lo que fluye.
(d) Verdadero.
\end{solution}

\begin{solution}{exo:g8:ice-on-earth:3}
(a) Un \omterm{def:g8:ice-on-earth:glacial-landforms}{valle en U}. (b) Un bloque errático. (c) Una \omterm{def:g8:ice-on-earth:ice-sheet}{plataforma de hielo}. (d) El \omterm{def:g8:ice-on-earth:permafrost}{permafrost}.
\end{solution}

\begin{solution}{exo:g8:ice-on-earth:4}
(a) La Antártida y Groenlandia. (b) La Antártida, unos \qty{58}{m}
(Groenlandia: unos \qty{7.4}{m}). (c) Alrededor de una décima parte (10~\%).
\end{solution}

\begin{solution}{exo:g8:ice-on-earth:5}
(a) Unos \num{7.05}~millones de~\unit{km^2} en 1979 y
\num{3.57}~millones de~\unit{km^2} en 2012. (b) $7.05 - 3.57 = 3.48$: unos
\num{3.5}~millones de~\unit{km^2}. (c) $3.48 \div 7.05 \approx 0.49$:
alrededor de la mitad. (d) Un año suelto depende del \omterm{def:g1:watching-the-weather:weather}{tiempo} de ese verano:
2012 queda muy por debajo de la línea. La línea de tendencia usa todos los
años, así que un verano excepcional apenas la modifica.
\end{solution}

\begin{solution}{exo:g8:ice-on-earth:6}
(a) $2004 - 1941 = 63$ años. (b) $12000 \div 63 \approx 190.5$: unos
\qty{190}{m} al año. (c) $2021 - 1984 = 37$ años; $3200 \div 37 \approx
86$: unos \qty{86}{m} al año. (d) El \omterm{def:g8:ice-on-earth:glacier}{glaciar} Muir, más del doble de
deprisa.
\end{solution}

\begin{solution}{exo:g8:ice-on-earth:7}
En la \emph{zona de acumulación}, en lo alto de la \omterm{def:g1:landscapes-around-us:landform}{montaña}, cae cada año
más \omterm{def:g4:clouds-rain-snow:precipitation}{nieve} de la que se funde. La \omterm{def:g4:clouds-rain-snow:precipitation}{nieve} nueva entierra y comprime la \omterm{def:g4:clouds-rain-snow:precipitation}{nieve}
vieja, que se transforma en \emph{neviza} al cabo de alrededor de un año y
después, más abajo y mucho más tarde, en hielo compacto. Por su propio peso,
el hielo \emph{fluye} lentamente \omterm{def:g1:landscapes-around-us:landform}{valle} abajo. Más abajo, en la \emph{zona
de ablación}, se funde más deprisa de lo que cae la \omterm{def:g4:clouds-rain-snow:precipitation}{nieve}; el hielo que
llega al frente es hielo que se formó muy arriba, a menudo hace mucho
tiempo.
\end{solution}

\begin{solution}{exo:g8:ice-on-earth:8}
(a) En el vaso B: el hielo fundido cae al agua y hace subir el nivel; en el
vaso A el nivel se queda donde estaba. (b) El vaso A es el \omterm{def:g8:ice-on-earth:ice-sheet}{iceberg} (hielo
que ya flota en el \omterm{def:g2:water-on-earth:ocean}{mar}); el vaso B es el \omterm{def:g8:ice-on-earth:ice-sheet}{inlandsis} (hielo sobre
tierra firme). (c) Una \omterm{def:g8:ice-on-earth:ice-sheet}{plataforma de hielo} retiene los \omterm{def:g8:ice-on-earth:glacier}{glaciares} que tiene
detrás. Cuando se rompe, esos \omterm{def:g8:ice-on-earth:glacier}{glaciares} se aceleran y vierten más hielo
continental en el \omterm{def:g2:water-on-earth:ocean}{mar}, y eso lo hace subir.
\end{solution}

\begin{solution}{exo:g8:ice-on-earth:9}
(a) Un \omterm{def:g8:ice-on-earth:glacial-landforms}{valle en U}, excavado por un \omterm{def:g8:ice-on-earth:glacier}{glaciar}. (b) Un \omterm{def:g8:ice-on-earth:glacier}{glaciar} lo transportó
desde muy arriba en el \omterm{def:g1:landscapes-around-us:landform}{valle}, donde hay granito, y lo dejó allí cuando el
hielo se fundió: es un bloque errático. (c) Una \omterm{def:g8:ice-on-earth:glacial-landforms}{morrena} lateral: su altura
muestra qué espesor llegó a tener el \omterm{def:g8:ice-on-earth:glacier}{glaciar}; el hielo llegó al menos hasta
esa altura en la ladera del \omterm{def:g1:landscapes-around-us:landform}{valle}.
\end{solution}

\begin{solution}{exo:g8:ice-on-earth:10}
(a) $7.4 \div 10 = \qty{0.74}{m}$. (b) $58 \div 100 = \qty{0.58}{m}$.
(c) $58 \div 7.4 \approx 7.8$: unas 8 veces más.
\end{solution}

\begin{solution}{exo:g8:ice-on-earth:11}
(a) $7.21 \times 10^6 = \qty{7210000}{km^2}$ y $4.65 \times 10^6 =
\qty{4650000}{km^2}$. (b) $7.21 - 4.65 = \num{2.56}\text{ millones de }\unit{km^2}$, es
decir, \qty{2560000}{km^2}. (c) $2.56 \div 3 \approx 0.85$: unos
\num{0.85}~millones de~\unit{km^2} por década, algo más que los
\num{0.76}~millones de~\unit{km^2} de la línea de tendencia; la media de
cada década conserva los altibajos de sus diez veranos, pero las dos
maneras de leer el gráfico muestran la misma gran pérdida. (d) La extensión
no puede bajar de cero, y una línea de tendencia solo describe los años que
se midieron: lo que ocurra después depende de cómo cambie el \omterm{def:g5:weather-and-climate:climate}{clima} del
Ártico, así que la línea no es una previsión.
\end{solution}

\begin{solution}{exo:g8:ice-on-earth:12}
(a) $120 \div 1 = 120$: la \omterm{def:g1:landscapes-around-us:water}{costa} estaba unos \qty{120}{km} más \omterm{def:g2:water-on-earth:ocean}{mar} adentro.
(b) Se había evaporado del \omterm{def:g2:water-on-earth:ocean}{océano}, había caído en forma de \omterm{def:g4:clouds-rain-snow:precipitation}{nieve} sobre los
\omterm{def:g3:continents-and-oceans:continent}{continentes} y se había quedado allí como hielo de los grandes
\omterm{def:g8:ice-on-earth:ice-sheet}{inlandsis}. (c) Una \omterm{def:g8:ice-on-earth:glacial-landforms}{morrena} frontal se forma en el extremo de un \omterm{def:g8:ice-on-earth:glacier}{glaciar}.
Una \omterm{def:g8:ice-on-earth:glacial-landforms}{morrena} en una \omterm{def:g1:landscapes-around-us:landform}{llanura} lejos de cualquier \omterm{def:g1:landscapes-around-us:landform}{montaña} muestra que un
\omterm{def:g8:ice-on-earth:ice-sheet}{inlandsis} llegó hasta esa \omterm{def:g1:landscapes-around-us:landform}{llanura}, en una época en que el hielo
cubría mucha más tierra que hoy.
\end{solution}

\begin{solution}{pb:g8:ice-on-earth:1}
\textbf{1.} $1080 - 0 = \qty{1080}{m}$.

\textbf{2.} 1960--1970: \qty{180}{m}; 1970--1980: \qty{150}{m};
1980--1990: \qty{190}{m}; 1990--2000: \qty{180}{m}; 2000--2010:
\qty{180}{m}; 2010--2020: \qty{200}{m}. El mínimo, en 1970--1980; el
máximo, en 2010--2020.

\textbf{3.} No: un \omterm{def:g8:ice-on-earth:glacier}{glaciar} en retroceso sigue fluyendo cuesta abajo. Su
frente retrocede porque el hielo se funde más deprisa de lo que llega.

\textbf{4.} Una \omterm{def:g8:ice-on-earth:glacial-landforms}{morrena} frontal: el hielo transportó derrubios de \omterm{def:g3:stones-pebbles-sand:rock}{roca} \omterm{def:g1:landscapes-around-us:landform}{valle}
abajo y los depositó en su frente, donde el hielo se fundía, mientras el
frente estuvo allí.

\textbf{5.} $1080 \div 6 = \qty{180}{m}$ por década.

\textbf{6.} $180 \div 10 = \qty{18}{m}$ al año.

\textbf{7.} $190 \div 18 \approx 10.6$: unas 11 veces más deprisa.

\textbf{8.} Menos: perdió más hielo en su zona de ablación del que ganó en
su zona de acumulación, así que su frente retrocedió.

\textbf{9.} En forma de U: un fondo ancho con laderas empinadas, porque el
\omterm{def:g8:ice-on-earth:glacier}{glaciar} ensanchó y profundizó el \omterm{def:g1:landscapes-around-us:landform}{valle} mientras lo llenaba.

\textbf{10.} $1980 - 1080 = \qty{900}{m}$.

\textbf{11.} $900 \div 18 = 50$ años: en \textbf{2070}.

\textbf{12.} La velocidad cambió de una década a otra (de 150 a
\qty{200}{m}), y el retroceso futuro depende de las nevadas y de las
temperaturas de verano de los próximos años: un ritmo pasado no es una
previsión.
\end{solution}
